%%%%%%%%%%%%%%%%%%%%%%%%%%%%%%%%%%%%%%%%%%%%%%%%%%%%%%%%%%%%%%%%%%%%%%%%
% Plantilla TFG/TFM
% Universidad de Murcia. Facultad de Informática
% Realizado por: José Manuel Requena Plens
% Modificado: Pablo José Rocamora Zamora
% Contacto: pablojoserocamora@gmail.com
%%%%%%%%%%%%%%%%%%%%%%%%%%%%%%%%%%%%%%%%%%%%%%%%%%%%%%%%%%%%%%%%%%%%%%%%


\chapter*{Preámbulo}
\thispagestyle{empty}

El avance de los LLMs ha transformado profundamente la manera en que accedemos y generamos información. Sin embargo, esta misma capacidad los convierte en una herramienta potencialmente peligrosa en un ecosistema mediático donde la desinformación se propaga a una velocidad que los mecanismos de verificación tradicionales son incapaces de absorber.

Este Trabajo Fin de Grado nace del pensamiento de que la tecnología que forma parte del problema puede también ser parte de la solución. A lo largo de estas páginas se describe el diseño, implementación y validación de un sistema de \textit{Fact-Checking} automatizado construido sobre una arquitectura de Generación Aumentada por Recuperación (RAG), con el objetivo de mitigar las alucinaciones de los modelos de lenguaje anclando sus respuestas en un corpus real de noticias.

\cleardoublepage %salta a nueva página impar

\chapter*{Agradecimientos}

\thispagestyle{empty}
\vspace{1cm}

En primer lugar, quiero expresar mi más sincero agradecimiento a mis tutores, cuya ayuda me ha ayudado a ir mucho más allá de lo que tenía previsto cuando empecé con este TFG. Sin su implicación, muchas de las decisiones técnicas que sustentan esta memoria no habrían madurado de la misma manera.

A mis padres, María José e Isidoro, no tengo palabras suficientes para agradeceros todo lo que habéis hecho por mí. Habéis sido el pilar sobre el que se ha sostenido no solo estos 4 años, sino cada etapa de mi vida. Gracias por haberme dado la oportunidad de estudiar lo que me gusta y por estar siempre ahí, sin condiciones y sin importar el momento.

A mi pareja Laura, gracias por aguantarme en los momentos de más estrés, hacerme ver lo mejor de mí y por no quejarte demasiado cuando he tenido la cabeza en otro sitio. Tu apoyo ha sido mucho más importante de lo que probablemente imaginas, y este trabajo al igual que cada logro que consigo también es un poco tuyo.

A mis abuelas y abuelos, por el cariño y el apoyo que han transmitido y transmiten con cada conversación, y por el orgullo que sienten de manera incondicional en todo lo que hago. Sois mi mayor ejemplo de perseverancia, solidaridad y amor. Por eso, y por todo lo que me habéis dado, gracias.

A mis amigos Antonio, Juanjo, Javi y Álvaro, que he conocido por el camino. En vosotros he encontrado a una familia con la que compartir no solo esta etapa, sino también muchas más que vendrán después. Gracias por los momentos de risas, las conversaciones y por hacer de la universidad un lugar que siempre recordaré.

Por último, a todas las personas que de una forma u otra han estado ahí durante este proceso: gracias. Saber que contaba con vosotros ha hecho que todo fuera mucho más llevadero.

\cleardoublepage %salta a nueva página impar

% Aquí va la dedicatoria si la hubiese. Si no, comentar la(s) linea(s) siguientes
%\chapter*{}
%\setlength{\leftmargin}{0.5\textwidth}
%\setlength{\parsep}{0cm}
%\addtolength{\topsep}{0.5cm}
%
%\begin{flushright}
%	\small\em{
%		A mis padres, María José e Isidoro, y a Laura.\\
%		Por todo lo que hay detrás de estas páginas.
%	}
%\end{flushright}


\cleardoublepage %salta a nueva página impar

% Aquí va la cita célebre si la hubiese. Si no, comentar la(s) linea(s) siguientes
\chapter*{}

\setlength{\leftmargin}{0.5\textwidth}
\setlength{\parsep}{0cm}
\addtolength{\topsep}{0.5cm}
\begin{flushright}
	\small\em{
		Buscando el bien de nuestros semejantes,\\
		encontramos el nuestro
	}
\end{flushright}

\begin{flushright}
	\small{
		Platon.
	}
\end{flushright}

\cleardoublepage %salta a nueva página impar

